\section{A literal quadratic-width, linear-depth theorem}
\label{sec:digital}

The logarithms in the preceding parameter-grid bound can also be removed
under a different structural restriction: after merging, the network behaves as a
collection of robust digital gates.  This assumption is natural for a
saturated implementation of a Boolean circuit, but is not claimed to hold
for typical gradient-trained analog networks.  The resulting prior still
has positive probability around an implementation of every eligible
Boolean circuit.  Its support definition does not use that circuit or its
width.

We retain the Gaussian tanh model, $H=L-1$, the readout energy
$E_v(w)=N\|v\|_2^2$, and the clipped score
$T(u)=\max\{-1/2,\min\{u,1/2\}\}$.
For balanced population partitions into $q$ groups at every hidden layer,
let $\bar w$ be the population-averaged full-width network.  Write its
exactly equivalent width-$q$ quotient as
\[
 h^0(x)=(x,1),\qquad a^s(x)=U^sh^{s-1}(x),\qquad
 h^s(x)=\tanh(a^s(x)),\qquad z_{\bar w}(x)=u^\top h^H(x).
\]
Thus $U^1$ has $d+1$ columns, subsequent $U^s$ have $q$ columns,
and $u$ has $q$ coordinates.

Fix, before observing labels,
\begin{equation}\label{eq:digital-calibration}
 \lambda=(d+L)\bigl(1+\log(d+2)\bigr),\qquad
 \Psi(q)=q\log^2(2+q),\qquad
 \eta=\frac1{16},\quad \kappa=\frac15.
\end{equation}
The event $\mathcal D_N$ consists of networks admitting such partitions,
with $1\le q\le N$, for which
\begin{align}
 \lambda\Psi(q)&\le E_v(w),\label{eq:digital-count}\\
 \sup_x|T(z_w(x))-T(z_{\bar w}(x))|&\le\frac1{16},
                    \label{eq:digital-merge}\\
 |h^s_a(x)|&\ge\frac{15}{16},
                    \label{eq:digital-saturation}\\
 |a^s_a(x)|&\ge\kappa\|U^s_{a,\cdot}\|_1,
 \qquad |z_{\bar w}(x)|\ge\kappa\|u\|_1.
                    \label{eq:digital-relative}
\end{align}
The last two lines hold for every Boolean input, every hidden layer,
and every indicated coordinate.  All rows and the readout in
\eqref{eq:digital-relative} are required to be nonzero.
The prior is the raw Gaussian conditioned on $\mathcal D_N$ and
\[
 \mathcal G_N=\{w: |z_w(x)|>1/2\text{ for every Boolean }x\}.
\]
Training conditions further on the observed noiseless signs.  We also
allow intersection with the original internal-noise event
$\mathcal R_{N,c}$, for any fixed $c>0$: clipped scores must change by at
most $1/16$ under every simultaneous post-activation and output disturbance
whose norm
\[
 \left(\sum_{s=1}^H\|\xi^s\|_2^2/N+|\xi^{\rm out}|^2\right)^{1/2}
\]
is at most $c/\sqrt{1+HE_v(w)}$.
This optional intersection does not change the conclusion below.

These are Borel restrictions involving finitely many inputs and partition
tuples at each width.  Saturation makes the populations nearly Boolean;
the relative row margin asks that the sign of each gate tolerate a fixed
fractional error in its aggregate incoming coefficients.  The stronger
population-count calibration in \eqref{eq:digital-count} is an additional
choice of prior: it compensates for the logarithmic gains used in the
Boolean implementation.  It is not implied by the original robustness
restriction.

\subsection{Exact budgets and the learning statement}

Let the Boolean target have layered width $C$, with $H$ stages and the
constant wires required by the NAND/COPY/NOT implementation counted.
Put
\begin{align}
 e_*&=\frac1{16\sqrt{2C}},\qquad \tau_0=e_*/4,\qquad
 b=\sqrt{16\log9+8\log4},\notag\\
 A&=\max\{32b,2\log(8/e_*)\},\qquad
 a=\sqrt\lambda\max\{4\log(2+C),256c\},\notag\\
 \rho&=\min\{1/4,A/[4(d+1)],Ae_*/(64\sqrt C)\},\notag\\
 q_0&=d+2+(H-1)C,\notag\\
 B_0&=9dA^2+18(H-1)CA^2+2a^2C,\notag\\
 D_{\rm dig}&=B_0+q_0\log\frac{e\sqrt{2\pi}}{2\rho}+(H-1)\log2.
 \label{eq:digital-target-budget}
\end{align}
For $0<\varepsilon,\delta<1$, define
\begin{equation}\label{eq:digital-stat-budget}
 \begin{aligned}
 B&=\max\{8,8[D_{\rm dig}+\log(6/\varepsilon)]\},\\
 Q&=\max\{q\in\mathbb N:q\ge1,\ \Psi(q)\le B/\lambda\},\\
 P_Q&=Q(d+2)+(H-1)Q^2.
 \end{aligned}
\end{equation}
The maximum exists and satisfies $Q\ge C$, as the proof verifies.

\begin{theorem}[Literal quadratic width and linear depth]
\label{thm:digital}
Under either prior just specified, all posterior normalizers are positive.
If
\begin{equation}\label{eq:digital-samples}
 n\ge\frac6\varepsilon
       \left[\log Q+P_Q\log33+\log(1/\delta)\right],
\end{equation}
then, with probability at least $1-\delta$ over iid training inputs,
the posterior gives the correct label probability strictly greater than
$2/3$ outside a set of test inputs of probability at most $\varepsilon$.
The same sample event works for every $N\ge C$ and every $0<t\le1$,
under an arbitrary common training and test distribution on the Boolean
cube.

There is an absolute constant $K$ such that the simpler sufficient bound
\begin{equation}\label{eq:digital-envelope}
 n\ge\frac{K(1+c^2)^2}{\varepsilon}
       \left[(d+2)R+LR^2+\log(1/\delta)\right],
 \qquad R=C+\frac{\log(6/\varepsilon)}{\lambda},
\end{equation}
implies the same conclusion.  In particular, at fixed input dimension,
calibration, accuracy, and confidence, the sufficient sample count is
$O(C^2L)$, with no hidden logarithm of $C$ or $L$.
\end{theorem}

The theorem covers the finite mean-field prior and every stronger scaling
$\gamma\ge1$, by taking $t(N)=N^{(1-\gamma)/2}$.  Its digital hypothesis is separate from ordinary internal-noise stability;
the conclusion also holds when both conditions are imposed. It strengthens
clean population averaging through the stated digital and count requirements.  It gives no claim about efficient posterior
sampling or typical training outcomes.

\subsection{A Gaussian cylinder inside the digital support}

\begin{lemma}[Digital target cylinder]\label{lem:digital-cylinder}
For every $N\ge C$ and $0<t\le1$, a set of raw probability at least
$\exp(-ND_{\rm dig}/t^2)$ belongs to
$\mathcal G_N\cap\mathcal D_N\cap\mathcal R_{N,c}$ and satisfies
$f(x)z_w(x)>1$ at every input.
\end{lemma}
\begin{proof}
Use the balanced cloned Boolean implementation with gate gain $A$ and
readout coefficient sum $a$.  In coordinates
$\sqrt d W^1$, $\sqrt N W^s$, and $Nv$, constrain first-layer coordinates,
every internal row's normalized clone-group directions, and all readout
coordinates to intervals of radius $\rho t$ around their teacher values.
The unconstrained internal blocks satisfy $E^sP=0$.  The same two-net
Gaussian estimate used for the target cylinder gives
\[
 \Pr\{\|E^s\|_{\rm op}>bt\}
 \le2\exp(2N\log9-Nb^2/8)=2\cdot4^{-N}\le\frac12.
\]
The blocks are independent.  Requiring all $H-1$ bulk bounds therefore
costs at most a factor $2^{-(H-1)}$.  This replaces a depth-dependent
operator-norm cutoff by a constant cutoff; the probability cost is
included in $D_{\rm dig}$.

The saturation induction for the cloned construction gives hidden RMS
error at most $e_*$ relative to its ideal Boolean state, even with RMS
noise at most $\tau_0$ separately at every layer and output noise at most
$\tau_0$.  Indeed its recurrence is
\[
 e_s\le 2e^{-A/2}+(8/A)\rho\sqrt C(1+\tau_0)
                 +(8b/A)e_{s-1}+\tau_0\le e_*.
\]
The first-layer perturbation initializes the same invariant, and the
signed readout is at least
$15a/16-\rho(1+\tau_0)-\tau_0>1$.
There are $Nq_0$ independent constrained scalar coordinates with sum of
squared centers at most $NB_0$.  The scalar Gaussian interval estimate
\[
 \Pr\{|Z-m|\le\rho t\}
 \ge\frac{2\rho}{e\sqrt{2\pi}}e^{-m^2/t^2}
 \quad (Z\sim\mathcal N(0,t^2))
\]
and $N/t^2\ge1$ give the claimed probability $e^{-ND_{\rm dig}/t^2}$.

Averaging within clone groups maps this cylinder into itself.  Both the
original and averaged networks therefore have clipped score $f/2$,
proving \eqref{eq:digital-merge} and the global margin.
Furthermore,
\[
 E_v(w)\ge a^2C/8\ge2\lambda C\log^2(2+C),
\]
which proves \eqref{eq:digital-count}.  The calibration $a\ge256c\sqrt\lambda$
also ensures that the original internal-noise ball lies in the preceding
$\tau_0$ product of balls, establishing the optional restriction.

It remains to check the digital conditions for the quotient.  Its hidden
coordinates are constant on groups of mass at least $1/(2C)$.  The RMS
error bound thus implies that every coordinate differs from its ideal
bit by at most $\sqrt{2C}e_*=1/16$, proving saturation.
Every ideal gate row has $\ell_1$ norm at most $3A$ and signed ideal
preactivation at least $A$.

In the first layer, the quotient row perturbation has $\ell_1$ norm at
most $(d+1)\rho t/\sqrt d\le A/4$.  Signed preactivation divided by row
norm is therefore at least $3/13>1/5$.
For an internal row, the perturbation coefficient from group $b$ is at
most $\rho t\sqrt{k_b/N}$.  Its total $\ell_1$ norm is at most
\[
 \rho\sqrt{2C}\le A/1024.
\]
Accounting also for the coordinatewise $1/16$ error of the preceding
state gives signed preactivation at least $831A/1024$ and row norm at
most $3073A/1024$.  Their ratio exceeds $1/5$.
Finally, the quotient readout perturbation has $\ell_1$ norm at most
$\rho$, so its signed output is at least $15a/16-\rho$ and its row norm
at most $a+\rho$.  As $a\ge4$ and $\rho\le1/4$, their ratio is at least
$14/17>1/5$.  All these rows are nonzero.  This verifies
\eqref{eq:digital-relative} and completes the support check.
\end{proof}

\subsection{A constant-precision digital code}

\begin{lemma}[Finite digital class]\label{lem:digital-code}
For each integer $Q\ge1$, the sign functions of networks in
$\mathcal G_N\cap\mathcal D_N$ admitting witnesses with $q\le Q$, over
all learner widths $N$, belong to a fixed finite class $\mathcal F_Q$
with
\[
 \log|\mathcal F_Q|\le\log Q+P_Q\log33.
\]
This class does not depend on the sample, learner width, or Gaussian scale.
\end{lemma}
\begin{proof}
The unit $\ell_1$ ball in $\mathbb R^m$ has an $\ell_1$ net of radius
$1/16$ with at most $33^m$ points.  To see this, take a maximal
$1/16$-separated set.  Disjoint balls of radius $1/32$ around its points
fit inside the ball of radius $1+1/32$; volume scaling gives the bound,
and maximality gives a net.

For a quotient satisfying the digital conditions, let
$b^s(x)=\operatorname{sign}(h^s(x))$.  Normalize each nonzero row by its
$\ell_1$ norm.  Since $\|h^s-b^s\|_\infty\le1/16$, replacing a preceding
hidden state by its bits changes a normalized preactivation by at most
$1/16$.  Each hidden gate and the final readout on those bits consequently
have the original sign with margin at least
\[
 \kappa-\eta=\frac{11}{80}.
\]
The first-layer input is already Boolean, including its constant
coordinate, so it has margin at least $\kappa$.

Replace every normalized row by a point in the fixed $1/16$ net for its
dimension, and use threshold gates instead of tanh.  On Boolean inputs
this changes each preactivation by at most $1/16<11/80$.  Induction through
the layers preserves every hidden sign and the readout sign.  At width
$q$, the product of the row-net sizes is at most $33^{P_q}$.  Summing over
$q\le Q$ gives at most $Q33^{P_Q}$ possibilities.  No count of adaptive
population partitions is required: every resulting quotient is already
covered by these same row nets.

Finally, the original global margin and
\eqref{eq:digital-merge} force the original and quotient scores to have
the same sign at every input.  Thus this code describes the actual
posterior sign functions, not just an approximate surrogate.
\end{proof}

The finite nets may use real coefficients.  The lemma counts a fixed
finite function class; it does not claim a rational bit encoder or an
efficient construction algorithm.

\subsection{Posterior learning and the explicit polynomial}

\begin{proof}[Proof of Theorem~\ref{thm:digital}]
Lemma~\ref{lem:digital-cylinder} lower-bounds the raw conditioning
probability by $e^{-ND_{\rm dig}/t^2}$ for every correctly labeled sample.  Under
the raw prior, $NE_v/t^2$ is chi-square with $N$ degrees of freedom.
For $B\ge8$, its exponential moment gives
\[
 \mu_{N,t}\{E_v>B\}\le e^{-NB/(8t^2)}.
\]
After conditioning on prior support and training labels, the mass of
this event is therefore at most
\[
 \exp\left[-\frac N{t^2}(B/8-D_{\rm dig})\right]\le\varepsilon/6.
\]
Every remaining predictor has a witness with $q\le Q$.  The definition
of $Q$ is legitimate because
$D_{\rm dig}\ge2a^2C\ge32\lambda\Psi(C)$, so $B/\lambda\ge256\Psi(C)$.

Every retained sign function lies in $\mathcal F_Q$ and is consistent
with the sample.  A fixed function of population error greater than
$\varepsilon/6$ is consistent with $n$ iid labels with probability at
most $e^{-n\varepsilon/6}$.  Lemma~\ref{lem:digital-code}, a union bound,
and \eqref{eq:digital-samples} imply that, with probability at least
$1-\delta$, every consistent member of $\mathcal F_Q$ has error at most
$\varepsilon/6$.  This event is independent of $N,t$.
On it, the full posterior Gibbs error is at most $\varepsilon/3$:
$\varepsilon/6$ for retained mass and at most $\varepsilon/6$ for
discarded mass.  The test mass where correct-label posterior probability
is at most $2/3$ is at most three times this Gibbs error.  This proves the
stated confidence guarantee, including arbitrary posterior mixtures.

For the simpler envelope, all constants below are absolute.  From
\eqref{eq:digital-target-budget},
\[
 A^2\le K_1\log^2(2+C),\qquad
 \rho^{-1}\le K_2(d+C+1),\qquad
 a^2\le K_3(1+c^2)\lambda\log^2(2+C).
\]
Using $d+LC\le(d+L)C$ and the dimension logarithm in
\eqref{eq:digital-calibration}, these inequalities give
\begin{equation}\label{eq:digital-cost-envelope}
 D_{\rm dig}\le K_4(1+c^2)\lambda\Psi(C).
\end{equation}
The interval logarithm contributes only
$O(\log^2(2+C)+\log(d+2))$, and the bulk cost is linear in $H$;
there is no logarithm of $L$ in this estimate.

For $k\ge1$, $\Psi(kC)\ge k\Psi(C)$.  Moreover, for $q\ge kC$,
\[
 \Psi(q)\ge k\Psi(C)+(q-kC)\log^2 3.
\]
Combining this with \eqref{eq:digital-cost-envelope} and the definition
of $B$ bounds
\[
 Q\le K_5\left((1+c^2)C+
                  \frac{\log(6/\varepsilon)}\lambda\right)
       \le K_5(1+c^2)R.
\]
Since $\log Q\le Q$ and
$P_Q\le(d+2)Q+LQ^2$, substitution in
\eqref{eq:digital-samples} proves \eqref{eq:digital-envelope} after
enlarging one absolute constant.
\end{proof}

\begin{corollary}[A polynomial accuracy envelope]
\label{cor:digital-polynomial}
For an absolute constant $K$, the sufficient sample condition
\begin{equation}\label{eq:digital-clean-envelope}
 n\ge\frac K\varepsilon\left[
 (1+c^2)^2(dC+LC^2)+\log^2(6/\varepsilon)+\log(1/\delta)\right]
\end{equation}
implies every conclusion of Theorem~\ref{thm:digital}.
\end{corollary}
\begin{proof}
Write $s=1+c^2$ and $v=\log(6/\varepsilon)>1$.  The preceding proof gives
$Q\le K_5(sC+v/\lambda)$.  Since $\lambda\ge d+L$,
\[
 (d+2)\frac v\lambda=O(v),\qquad
 L\frac{v^2}{\lambda^2}\le v^2,\qquad
 2LsC\frac v\lambda\le2sCv\le s^2C^2+v^2.
\]
Also $(d+2)sC=O(s^2dC)$, because $d\ge1$ and $s\ge1$.
Consequently
$\log Q+P_Q\log33=O(s^2(dC+LC^2)+v^2)$.
Insert this estimate into \eqref{eq:digital-samples}.
\end{proof}

\paragraph{What the sharpening does and does not establish.}
The $C^2L$ conclusion uses both digital relative margins and the stronger
readout-energy population calibration.  It does not assert the same bound
for the original internal-noise prior, for clean neuron averaging alone,
or for the unrestricted Gaussian prior.  Every eligible Boolean circuit
has a positive-mass witness, but that is an existence statement about an
implementation, not evidence that generic trained networks satisfy the
support condition.  The proof is a conventional argument; no Lean
formalization is claimed for this theorem.
